\section{Conclusion}

We asked whether a digital mind's account of itself can be taken as evidence
about it, and we answered by building the mind out of parts and intervening
between them. A commitment planted in the subject's reasoning steered the answer
while the mind's report about its own thinking rarely named it. Holding a window
fixed and changing only the model reading it moved the readout far more than
changing whose conversation was read, and the subject's own weights gave a reader
no advantage. Reading the subject without asking it anything separated held-out
scenarios that a probe's questions did not. We therefore contend that a
self-report about a digital mind should not be published without the calibration
that says how much of it belongs to the probe.
